% from ega4-8-fr.tex
\oldpage[IV-3]{54}
Reste à prouver que (8.14.2.4) est surjective. Soit donc $v:Z\to V$ un
$S$-morphisme~; par hypothèse il existe un $\lambda$ et un $S$-morphisme
$w_\lambda:Z_\lambda\to X$ tels que $j\circ v$ se factorise en
\[
Z\xrightarrow{u_\lambda}Z_\lambda\xrightarrow{w_\lambda}X,
\]
et tout revient à prouver qu'il existe $\mu\geqslant\lambda$ tel que le
morphisme
\[
Z_\mu\xrightarrow{u_{\lambda\mu}}Z_\lambda
\xrightarrow{w_\lambda}X
\]
(où $u_{\lambda\mu}$ est le morphisme de transition) se factorise en
\[
Z_\mu\xrightarrow{v_\mu}V\xrightarrow{j}X.
\]
Posons, pour tout $\mu\geqslant\lambda$,
$U_\mu=u_{\lambda\mu}^{-1}(w_\lambda^{-1}(V))$. On a
\[
u_\mu^{-1}(U_\mu)=u_\lambda^{-1}(U_\lambda)
=u_\lambda^{-1}(w_\lambda^{-1}(V))=v^{-1}(V)=Z.
\]
Comme les $Z_\lambda$ sont quasi-compacts et que les $U_\mu$, étant ouverts,
sont ind-constructibles (1.9.6), on déduit de (8.3.4) qu'il existe
$\mu\geqslant\lambda$ tel que $U_\mu=Z_\mu$. C.Q.F.D.

\begin{remark}[8.14.3]
\label{IV.8.14.3-fr}
Le fait que l'application (8.14.2.1) soit injective lorsque $f$ est localement
de type fini (8.8.2, (i)) amène naturellement à se demander si ce résultat
admet aussi une réciproque. Il n'en est rien, même lorsque $S$ et $X$ sont
affines, car il existe des monomorphismes $X\to S$ qui ne sont pas de type fini
(I, 2.4.2), et qui mettent donc en défaut cette conjecture.
\end{remark}

% (continued in ega4-9-fr.tex, before its first page marker)
\section{Propriétés constructibles}
\label{section:IV.9-fr}

Soient $S$ un préschéma, $f:X\to S$ un morphisme de
\emph{présentation finie} (1.6.1), $\mathcal F$ un $\mathcal O_X$-Module
quasi-cohérent de \emph{présentation finie}. On se propose dans ce paragraphe
de donner des critères assurant, par exemple, que l'ensemble des $s\in S$ tels
que le $k(s)$-préschéma
$X_s=f^{-1}(s)=X\times_S\operatorname{Spec}(k(s))$ a une certaine propriété,
ou tels que le $\mathcal O_{X_s}$-Module
$\mathcal F_s=\mathcal F\otimes_{\mathcal O_X}k(s)$ a une certaine propriété,
est \emph{localement constructible}, ou tout au moins
\emph{ind-constructible} (1.9.4)~; on verra que c'est le cas pour la plupart
des propriétés qui s'introduisent en Géométrie algébrique. En supposant
seulement $f$ \emph{localement de présentation finie}, on donnera aussi (9.9)
des critères pour que l'ensemble des points $x\in X$ où la fibre $X_{f(x)}$
(ou le $\mathcal O_{X_{f(x)}}$-Module $\mathcal F_{f(x)}$) a une certaine
propriété, soit \emph{localement constructible}. On verra au \S~12 que ces
résultats, joints à l'hypothèse supplémentaire que $f$ est \emph{plat}
(resp. \emph{propre et plat}), permettent de prouver que les ensembles
considérés dans $X$ (resp. dans $S$) sont même \emph{ouverts} dans de nombreux
cas.

\subsection{Le principe de l'extension finie}
\label{subsection:IV.9.1-fr}

\begin{proposition}[9.1.1]
\label{IV.9.1.1-fr}
\textnormal{(Principe de l'extension finie).} Soient $k$ un corps,
$\mathcal C$ un ensemble d'extensions de $k$. On suppose vérifiées les
conditions suivantes~:
\begin{enumerate}[label=(\roman*)]
  \item Si $K\in\mathcal C$ et s'il existe un $k$-homomorphisme $K\to K'$
  (où $K'$ appartient à l'univers où l'on se place), alors
  $K'\in\mathcal C$.
  \item Si $K\in\mathcal C$, il existe une sous-extension $K'$ de $K$, de
  type fini sur $k$, telle que $K'\in\mathcal C$.
  \item Si $K\in\mathcal C$ est le corps de fractions d'une $k$-algèbre
  intègre de type fini $A$, il existe $f\in A-\{0\}$ tel que pour tout idéal
  maximal $\mathfrak m$ de $A$ on ait $A_f/\mathfrak m\in\mathcal C$.
\end{enumerate}
